\subsection{A common maximal-rank incidence for odd principal-cube systems}

Work over an algebraically closed field \(k\) of characteristic two.
All Cartier maps in this subsection are relative.  The geometric
Moore identification and its universal singular-parameter theorem
are used in their separately proved forms.  We keep the multiplication
map, its dual contractions, and the actual coherent Frobenius distinct.

\begin{lemma}[The other contraction of the same Moore tensor]
\label{lem:DE-common-incidence-multiplication}
\status{Candidate; complete typed proof supplied.}
Let \((B,\Theta)\) be a principally polarized abelian surface,
put \(M=\mathcal O_B(3\Theta)\), and let \(P\) be a nontrivial
line bundle of odd order in \(\operatorname{Pic}^0(B)(k)\).
Let \(Q\) be its unique half in the odd cyclic subgroup it
generates, so \(Q^2=P\).  For a general \(s\in H^0(B,M)\),
the relative Cartier multiplication map
\[
 H^0(B,MP)\longrightarrow
 H^0(B^{(2)},M^{(2)}Q^{(2)}),\qquad
 t\longmapsto C_{\rm rel}(st\omega_B)/\omega_B'
                                                        \tag{DE-CI1}
\]
is an isomorphism.  Both spaces have dimension nine.
\end{lemma}
\begin{proof}
First consider the balanced pairing between \(MQ^{-1}\) and
\(MQ\), with target \(M^{(2)}\).  Proposition
\ref{prop:ppav-Cartier-cross-Moore-identification} gives bases,
indexed by \(H=\mathbf F_3^2\), in which its tensor is
\[
 \beta(e_i,f_j)=c_{(i-j)/2}\,g_{(i+j)/2},              \tag{DE-CI2}
\]
up to a single nonzero scalar.  Here the parameter \(c\)
belongs to the theta evaluation at the corresponding translation.
It is not an identically singular Moore parameter, by
Theorem~\ref{thm:ppav-non-two-torsion-cross-Cartier}, since
\((Q^{-1})^2\ne\mathcal O_B\).

If \(v=\sum_i v_i e_i\), the matrix of \(\beta(v,-)\),
with output row \(k\) and input column \(j\), is
\[
             c_{k-j}v_{2k-j}=c_{k-j}v_{-(k+j)}.
                                                        \tag{DE-CI3}
\]
The last equality is in the index group \(H\), where \(2=-1\).
Thus it is exactly the same Moore pencil with independent
variables \(x_h=v_{-h}\).  Its determinant is a nonzero
polynomial in \(v\).  This is an identity of the full tensor;
generic rank of one contraction alone is not being transferred
to a different contraction without that identity.

Choose \(u\in B(k)\) with \(\phi_M(u)=Q\).  Translation
by \(-u\) identifies the two inputs of (DE-CI1) with
\(MQ^{-1}\) and \(MQ\), respectively.  More explicitly,
\(t_{-u}^*M\simeq MQ^{-1}\), and translation preserves the
isomorphism class of every line bundle in \(\operatorname{Pic}^0\).
The target of (DE-CI1) is correctly typed because
\[
 F_B^*(M^{(2)}Q^{(2)})=M^2Q^2=M^2P.
\]
Translation of that target by \(-F_B(u)\) identifies it with
\(M^{(2)}\).  Naturality of relative Cartier under translation,
with invariant volumes, therefore identifies (DE-CI1) with
the balanced multiplication map.  Choices of line isomorphisms
only multiply the comparison by nonzero constants.  The input
translation is a fixed linear isomorphism, so the nonzero
determinant polynomial remains nonzero on \(H^0(B,M)\).
\end{proof}

\begin{proposition}[One section and its actual left kernel attain both maxima]
\label{prop:DE-common-Cartier-incidence}
\status{Candidate; complete geometric proof supplied.}
Let \(B\) be a nonordinary principally polarized abelian surface,
let \(M=\mathcal O_B(3\Theta)\), and let
\(\pi:A\to B\) be an odd-degree etale isogeny of degree \(n\).
Put \(L=\pi^*M\), \(V=H^0(A,L)\), \(W=H^0(A^{(2)},L^{(2)})\),
and \(h=\dim V=\dim W=9n\).  Define
\[
 \beta:\Lambda^2V\longrightarrow W,\qquad
 \beta(s,t)=C_{\rm rel}(st\omega_A)/\omega_A',
 \qquad M_s=\beta(s,-).
                                                        \tag{DE-CI4}
\]
There is a nonempty open subset of \(H^0(B,M)\) such that,
for every section \(s_0\) in it and \(s=\pi^*s_0\),
the following statements hold:
\begin{enumerate}
\item the divisors \(Y_0=(s_0=0)\) and \(Y=(s=0)\) are smooth
and connected, and the actual relative coherent Frobenius of \(Y\)
is bijective;
\item \(M_s\) has rank \(h-1\), with right kernel \(ks\);
\item if \(\lambda_s\in W^*\) spans its actual left kernel,
then the alternating form
\(B_{\lambda_s}(v,w)=\lambda_s\beta(v,w)\) has rank \(h-1\).
\end{enumerate}
In particular the section-rank open and the maximal dual-rank
open occur on the same incidence component, at an actual common pair.
\end{proposition}
\begin{proof}
The character decomposition of the odd etale isogeny is
\[
 \pi_*\mathcal O_A=\bigoplus_\chi P_\chi,
 \qquad P_1=\mathcal O_B,
 \qquad V=\bigoplus_\chi H^0(B,MP_\chi).
                                                        \tag{DE-CI5}
\]
For multiplication by a trivial-character section \(s_0\),
Cartier sends the \(P_\chi\)-summand to the target summand
\(M^{(2)}Q_\chi^{(2)}\), where \(Q_\chi^2=P_\chi\)
in the odd character group.  Thus the nontrivial character
summands are permuted, with no mixing, by taking their unique
halves.  Lemma~\ref{lem:DE-common-incidence-multiplication}
makes every one of these nine-by-nine blocks invertible on
a nonempty open subset of the same section space \(H^0(B,M)\).

On the trivial summand choose the theta basis and target
identification of Theorem~\ref{thm:ppav-level3-Cartier-symmetric}.
Its fully symmetric trilinear tensor makes the multiplication
matrix symmetric and alternating.  On its rank-eight open,
its right and left kernel coordinate vectors are both \(s_0\).
Equivalently, there is a fixed linear isomorphism
\[
 J:H^0(B,M)\xrightarrow{\ \sim\ }
            H^0(B^{(2)},M^{(2)})^*
\]
such that the trivial-block left kernel is \(J(s_0)\).
The target here is dual to the Frobenius-twisted section space;
\(J\) is the finite-theta representation identification, not
an untyped identification of the variety with its twist.

Extend \(J(s_0)\) by zero on every nontrivial target character
and call the resulting functional \(\lambda_s\).  After the
nontrivial multiplication blocks are made invertible, this is
the actual unique left kernel of \(M_s\).  That map then has
rank \(8+9(n-1)=h-1\), and alternation identifies its right
kernel as \(ks\).

It remains to impose the dual rank without changing that left
kernel.  A functional supported on the trivial target character
pairs only the trivial summand and the inverse pairs
\((\chi,\chi^{-1})\).  Each nontrivial cross determinant is
a nonzero polynomial on
\(H^0(B^{(2)},M^{(2)})^*\), by
Theorem~\ref{thm:ppav-non-two-torsion-cross-Cartier}.
Composing with the fixed linear isomorphism \(J\) makes it a
nonzero polynomial on the same nine-dimensional section space.
The fully symmetric trivial tensor identifies its dual contraction
at \(J(s_0)\) with the multiplication contraction at \(s_0\),
so that block has rank eight as well.  Intersect all these
finitely many determinant opens with the multiplication opens,
the degree-one rank-eight open, and the degree-one smooth/relative-
Frobenius-bijective open from Proposition~\ref{prop:DE-A3-ppav-degree-one}.
They are nonempty opens in one integral affine space over the
infinite field \(k\), so their intersection is nonempty.
At every point of the intersection,
\[
                 \operatorname{rank}B_{\lambda_s}
                      =8+18(n-1)/2=h-1.
                                                        \tag{DE-CI6}
\]

For completeness, the same choice gives the asserted actual
coherent conclusion.  The etale pullback of \(Y_0\) is smooth.
It is a divisor of the ample line bundle \(L\) on an abelian
surface, and \(H^1(A,L^{-1})=0\) by Serre duality and ample-line
vanishing; hence \(H^0(Y,\mathcal O_Y)=k\), so it is connected.
In the dual Cartier description of coherent Frobenius, all
nontrivial character blocks lie in the relative subspace and
are permuted bijectively.  Every block is invertible, so every
finite character cycle is semilinearly bijective.  The remaining
trivial relative block is precisely the degree-one relative
Frobenius for \(Y_0\), already chosen bijective.  Their direct
sum is the actual relative Frobenius for \(Y\).  This uses the
actual character decomposition; it does not assert that an
arbitrary etale pullback of a sharp curve remains sharp.
\end{proof}

\begin{corollary}[The dual kernel map is birational for every fixed odd degree]
\label{cor:DE-odd-principal-cube-kernel-birational}
\status{Candidate; complete deduction supplied.}
For every fixed nonordinary principally polarized abelian surface
\((A,\Theta)\) and every odd \(d\ge1\), set
\(L=\mathcal O_A(3d\Theta)\), \(h=9d^2\), and define \(\beta\)
as in (DE-CI4).  On the nonempty rank-\(h-1\) open
\(\mathcal U\subset\mathbf P(H^0(A,L))\), the actual left-kernel map
\[
 \kappa:\mathcal U\longrightarrow
        \mathbf P\bigl(H^0(A^{(2)},L^{(2)})^*\bigr)
\]
is separably dominant and birational.
\end{corollary}
\begin{proof}
Use the outgoing maximal-isotropic quotient in
Corollary~\ref{cor:fixed-ppav-all-odd-maximal-dual-Cartier}:
\(A\to B\) is etale of degree \(d^2\), and
\(L\) is the pullback of a cubic symmetric principal line bundle
on the nonordinary target \(B\).  The preceding proposition
supplies an actual point \([s]\in\mathcal U\) at which
\(\operatorname{rank}B_{\kappa(s)}=h-1\).
Proposition~\ref{prop:root-dual-kernel-differential} now gives
\(\operatorname{rank}d\kappa_{[s]}=h-1\), so \(\kappa\)
is separably dominant.  On the dense open where \(B_\lambda\)
has rank \(h-1\), its unique radical line recovers \([s]\).
This is the rational inverse, exactly as in
Corollary~\ref{cor:root-dual-kernel-birational}.
\end{proof}

\begin{remark}[What the common incidence does not prove]
\label{rem:DE-common-Cartier-incidence-scope}
\status{Candidate scope; no full monodromy or divisorial-inertia conclusion.}
The proof produces a common point in the actual incidence, rather
than merely separate generic-rank statements for \(M_s\) and
\(B_\lambda\).  It also proves generic sharpness for this particular
pullback family of cubic theta curves.  It does not establish that
every corank-one Hasse divisor meets the maximal-differential open,
that its actual local inertia is an entire hyperplane radical,
or that these radicals normally generate the global group.
Those are separate geometric requirements.  No new full-Lang or
monodromy degree is claimed by this subsection.

Every new tensor identity here is the explicit index equality
(DE-CI3), and every new geometric comparison uses etale pullback
or a specified relative-Frobenius target.  No additional numerical
experiment is needed or claimed; the complete universal Moore
certificate and the prior full theta-source audits are separate
retained dependencies.
\end{remark}
